% =========================================================
% What Does Co-Cauchy Mean?
% Source: volume-iii/analysis/metric-spaces/notes/completeness/
%
% Purpose: Give intuitive and formal explanation of the
%          co-Cauchy equivalence relation used in the completion.
% Dependencies: def:cauchy-sequence, def:complete-metric-space
% =========================================================

\subsubsection{Co-Cauchy Sequences: Two Sequences Heading for the Same Missing Point}
\label{sec:co-cauchy-intuition}

The completion construction takes an incomplete metric space $(X, d)$
and ``fills in'' every missing limit.
The key idea is to identify two Cauchy sequences that are heading
toward the \emph{same} missing point, even if that point doesn't exist in $X$ yet.
This is the notion of \emph{co-Cauchy sequences}.

\begin{remark*}[The core intuition]
\label{rem:co-cauchy-intuition}
Imagine you are in $\mathbb{Q}$ and you have two sequences:
\begin{align*}
a_n &= 1.4,\; 1.41,\; 1.414,\; 1.4142,\; \ldots \\
b_n &= 1.41,\; 1.414,\; 1.4142,\; 1.41421,\; \ldots
\end{align*}
Both are Cauchy sequences in $\mathbb{Q}$.
Both are converging toward $\sqrt{2}$, which does not exist in $\mathbb{Q}$.

Are $(a_n)$ and $(b_n)$ describing the same missing point, or different ones?
Clearly the same one: $\sqrt{2}$.

How do we know they are heading for the same point, \emph{without mentioning} $\sqrt{2}$?
By checking whether $d(a_n, b_n) \to 0$ --- the distance between corresponding
terms goes to zero.
\end{remark*}

\begin{tcolorbox}[colback=propbox, colframe=propborder, arc=2pt,
  left=6pt, right=6pt, top=4pt, bottom=4pt,
  title={\small\textbf{Definition (Co-Cauchy Sequences)}},
  fonttitle=\small\bfseries]
\begin{definition}[Co-Cauchy Sequences]
\label{def:co-cauchy}
Let $(X,d)$ be a metric space.
Two Cauchy sequences $(p_n)$ and $(q_n)$ in $X$ are called
\emph{co-Cauchy} (or \emph{equivalent}) if
\[
d(p_n, q_n) \;\to\; 0
\quad\text{as } n \to \infty.
\]
We write $(p_n) \sim (q_n)$.
\end{definition}
\end{tcolorbox}

\begin{remark*}[Fully quantified form]
\[
(p_n) \sim (q_n)
\;\iff\;
\forall \varepsilon > 0,\; \exists N \in \mathbb{N},\; \forall n \ge N,\; d(p_n, q_n) < \varepsilon.
\]
\end{remark*}

\begin{proposition}[Co-Cauchy is an equivalence relation]
\label{prop:co-cauchy-equivalence-relation}
The relation $\sim$ on the set of all Cauchy sequences in $(X,d)$ is an equivalence relation:
\begin{enumerate}[label=(\roman*)]
  \item \textbf{Reflexive:} $(p_n) \sim (p_n)$, since $d(p_n, p_n) = 0 \to 0$.
  \item \textbf{Symmetric:} If $(p_n) \sim (q_n)$ then $d(p_n, q_n) \to 0$, and since
        $d(q_n, p_n) = d(p_n, q_n)$ by symmetry, $(q_n) \sim (p_n)$.
  \item \textbf{Transitive:} If $(p_n) \sim (q_n)$ and $(q_n) \sim (r_n)$, then
        $d(p_n, r_n) \le d(p_n, q_n) + d(q_n, r_n) \to 0 + 0 = 0$, so $(p_n) \sim (r_n)$.
\end{enumerate}
\end{proposition}

\begin{example}[Co-Cauchy in $\mathbb{Q}$: both sequences heading for $\sqrt{2}$]
\label{ex:co-cauchy-sqrt2}
In $\mathbb{Q}$, let
\begin{align*}
a_n &= \text{decimal truncation of } \sqrt{2} \text{ to } n \text{ places}, \\
b_n &= a_n + \tfrac{1}{10^n}.
\end{align*}
Both $(a_n)$ and $(b_n)$ are Cauchy in $\mathbb{Q}$.
Moreover,
\[
|a_n - b_n| = \tfrac{1}{10^n} \to 0,
\]
so $(a_n) \sim (b_n)$.
They are co-Cauchy: both represent the same missing point $\sqrt{2}$.
\end{example}

\begin{example}[Not co-Cauchy: heading for different missing points]
\label{ex:not-co-cauchy}
In $\mathbb{Q}$, let $(a_n)$ be the decimal truncations of $\sqrt{2}$
and $(b_n)$ be the decimal truncations of $\pi$.
Both are Cauchy in $\mathbb{Q}$, but
\[
|a_n - b_n| \to |\sqrt{2} - \pi| \approx 1.72 \neq 0.
\]
So $(a_n) \not\sim (b_n)$: they are heading for different missing points.
\end{example}

\begin{remark*}[How the completion uses co-Cauchy sequences]
The completion $\widetilde{X}$ is defined as the set of all equivalence classes
$[(p_n)]$ of Cauchy sequences under $\sim$.
Each equivalence class represents one ``missing point.''

\begin{itemize}
  \item If $x \in X$, the constant sequence $(x, x, x, \ldots)$ is Cauchy, and
        its class $[(x,x,x,\ldots)]$ represents $x$ itself.
        This is how $X$ embeds in $\widetilde{X}$.
  \item If $(p_n)$ is Cauchy in $X$ with no limit in $X$, then
        $[(p_n)]$ is a new point in $\widetilde{X}$: the missing limit.
\end{itemize}

The metric on $\widetilde{X}$ is defined by
\[
\widetilde{d}\bigl([(p_n)], [(q_n)]\bigr) := \lim_{n \to \infty} d(p_n, q_n).
\]
This limit exists because the sequence $d(p_n, q_n)$ is Cauchy in $\mathbb{R}$
(a fact proved using the four-point inequality --- see Theorem~\ref{thm:completion}),
and $\mathbb{R}$ is complete.
\end{remark*}

\begin{remark*}[Why this is natural]
The completion of $\mathbb{Q}$ under this construction gives $\mathbb{R}$.
This is an alternative to the Dedekind cut construction:
instead of defining real numbers as sets of rationals,
we define them as equivalence classes of rational Cauchy sequences.
The number $\sqrt{2}$, for instance, is represented by the class
$[(1.4, 1.41, 1.414, \ldots)]$.
All Cauchy sequences converging to $\sqrt{2}$ fall in the same class.
\end{remark*}
